\documentclass[12pt]{article}

\usepackage[margin=1.18in]{geometry}
\usepackage{amsmath,amssymb,amsthm,mathtools}
\usepackage{array}
\usepackage{booktabs}
\usepackage{longtable}
\usepackage{hyperref}

\newtheorem{theorem}{Theorem}
\newtheorem{corollary}{Corollary}
\newtheorem{proposition}{Proposition}
\theoremstyle{definition}
\newtheorem{definition}{Definition}
\theoremstyle{remark}
\newtheorem{remark}{Remark}

\newcommand{\A}{\mathcal A}
\newcommand{\B}{\mathcal B}
\newcommand{\C}{\mathbb C}
\newcommand{\D}{\mathcal D}
\newcommand{\M}{\mathcal M}
\newcommand{\N}{\mathcal N}
\newcommand{\R}{\mathbb R}
\newcommand{\Z}{\mathcal Z}
\setlength{\emergencystretch}{3em}
\hypersetup{
  pdftitle={Finite Central Invisibility in Operational Reconstructions of Tensor Factorization and Horizon Entropy},
  pdfsubject={Operator algebras, algebraic quantum field theory, and horizon entropy},
  pdfkeywords={AQFT, von Neumann algebras, finite centers, tensor factorization, crossed products, DHR sectors, horizon entropy}
}

\title{Finite Central Invisibility in Operational Reconstructions of Tensor
Factorization and Horizon Entropy}
\author{}
\date{June 12, 2026}

\begin{document}
\maketitle

\begin{abstract}
Operational reconstructions in algebraic quantum field theory and quantum
information normally reconstruct the object that is supplied: a represented von
Neumann algebra, a local net, a sector category, a channel tested by specified
preparations and effects, or a Type II algebra with a trace normalization.  This
note proves a finite-central converse.  On comparison classes closed under
finite central direct sums, finite central refinements, and shared finite
centers in proposed bipartitions, represented or retracted operational data
determine full central support, central cardinality, tensor factorization,
central channel dynamics, retained sector labels, and absolute finite-center
entropy contributions exactly when the matching separating or calibrating
hypothesis is supplied.  If it is not supplied, explicit finite-central
same-data pairs change the target invariant.

The result is a no-go theorem, not a new positive reconstruction of quantum
gravity.  It says that normal state/effect data factored through a quotient do
not decide whether a nonzero finite central summand has been omitted; generated
algebra correlations do not reconstruct an unquotiented spatial tensor product
when multiplication has a kernel; and relative-entropy, recovery, modular-core,
or crossed-product entropy data do not fix an absolute horizon entropy
coefficient while a same-data finite-center or affine/trace residual has
nonzero area density.  Standard positive results--Tomita--Takesaki theory,
Osterwalder--Schrader reconstruction, DHR/Doplicher--Roberts reconstruction,
the split property, Fewster--Verch dynamical locality, process tomography,
JLMS/OAQEC, and gravitational crossed-product entropy--are used as inputs and
are left intact.  A Lean 4 file formalizes the finite logical descent kernels;
the analytic operator-algebraic reductions are the standard cited hypotheses.
\end{abstract}

\paragraph{Keywords.}
Algebraic quantum field theory; von Neumann algebras; finite centers; locally
covariant quantum field theory; DHR sectors; split property; process
tomography; modular crossed products; black-hole entropy.

\paragraph{MSC 2020.}
Primary 81T05, 46L10, 46L55; secondary 81P45, 83C57.

\section{Question and Contribution}

Many reconstruction arguments in quantum field theory and quantum information
start from operational data: normal state values, effects, channels, relative
entropies, modular flows, sector categories, or local probe readouts.  The
positive theorems then reconstruct the algebraic or geometric structure that
was actually supplied.  Tomita--Takesaki theory acts on a specified represented
von Neumann algebra and weight.  Osterwalder--Schrader reconstruction starts
from reflection-positive Euclidean data after quotienting by the null space.
DHR and Doplicher--Roberts reconstruction start from a specified sector
category of a specified observable net.  Locally covariant QFT and
Fewster--Verch dynamical locality are formulated for a supplied functor and its
relative Cauchy evolution.  Process tomography is relative to specified
preparations and effects; no-broadcasting constrains physical channels for
noncommuting state families, not unsupplied commutative central coordinates.
JLMS/OAQEC and crossed-product entropy use specified code algebras, represented
algebras, traces, and normalizations
\cite{TakesakiI,TakesakiII,OS1,OS2,DHR1,DHR2,DRDuality,DR,BFV,FewsterVerch,
FewsterVerchMeasurements,Choi,Jamiolkowski,NoBroadcasting,Petz,JLMS,HarlowRT,DHW,
WittenCrossed,LongoWitten,CLPW,CPW}.

The question here is the converse finite-center question.  Suppose the data
have already passed through a represented quotient, a central-state retraction,
a diagonal restriction, a generated algebra, a restricted input/effect family,
or an entropy calibration that does not see all finite central directions.  Can
the data nevertheless reconstruct the full ambient central support, the number
of finite central atoms, the unquotiented tensor product, or the absolute
finite-center contribution to a horizon entropy coefficient?

The answer is no in the precise fiber sense below.  A target invariant is
reconstructible from a labelled data map exactly when it is constant on the
data fibers.  The theorem identifies the finite-central fibers for the standard
operator-algebraic data maps named above.  It also records the exact positive
exits.  Full support requires central support faithfulness.  Central atoms and
weights require central separation.  Spatial tensor product coordinates require
tested-kernel separation, and with full generated-algebra tomography this is
injectivity of multiplication, supplied in AQFT by split or
W*-tensor-product independence.  Finite central process dynamics require
input spanning and output separation.  Central sector labels require the
supplied sector datum to retain those labels.  Absolute entropy and horizon
coefficients require trace, entropy, affine area, and geometric area-scale
calibrations that kill same-data residuals.

\paragraph{Novelty claim.}
No individual ingredient in this note is advertised as new.  Weak-* closed
ideals of von Neumann algebras are central summands; normal representations
have central support; normal functionals separate von Neumann algebras when
enough of them are supplied; split inclusions give tensor-product
independence; DHR/DR reconstruct from a supplied sector category; and Type II
entropy constructions are relative to a specified algebra and trace.  The
claim is the joint finite-central descent criterion: for the full-algebra
targets listed below, represented, retracted, restricted, generated, or tested
operational data determine the target if and only if the corresponding
finite-central separator, tensor-kernel separator, sector-character datum,
quotient/exclusion rule, or entropy/area calibration is supplied.  If an
existing paper already states this joint iff theorem, this note should be read
as an application to operational spacetime and horizon-entropy reconstruction
rather than as a new theorem.

\section{Data Descent}

\begin{definition}[Reconstruction from labelled data]
Let \(\mathcal X\) be a comparison class, \(D:\mathcal X\to\D\) a labelled data
map, and \(I:\mathcal X\to Y\) a target invariant.  The target is
reconstructible from \(D\) on \(\mathcal X\) if there is a map
\(R:D(\mathcal X)\to Y\) such that \(I=R\circ D\).
\end{definition}

\begin{proposition}[Fiber criterion]
\label{prop:fiber}
A target \(I:\mathcal X\to Y\) is reconstructible from \(D:\mathcal X\to\D\)
if and only if \(I\) is constant on every fiber of \(D\).  Equivalently, no
pair \(X,Y\in\mathcal X\) has \(D(X)=D(Y)\) but \(I(X)\ne I(Y)\).
\end{proposition}

\begin{proof}
If \(I=R\circ D\), then \(D(X)=D(Y)\) implies \(I(X)=R(D(X))=R(D(Y))=I(Y)\).
Conversely, if \(I\) is constant on data fibers, define
\(R(d)=I(X)\) for any \(X\) with \(D(X)=d\).  Fiber constancy makes this
well-defined.
\end{proof}

This elementary criterion is the only logical engine in the paper.  The
operator-algebraic work is to identify when ordinary operational data maps have
finite central fibers on which a full-algebra target changes.

\section{Main Finite-Central Descent Theorem}

\begin{theorem}[Finite-central operational descent]
\label{thm:main}
Let \(\mathcal X\) be a comparison class of von Neumann-algebraic local-net,
locally covariant, code-algebra, or finite process models.  Let
\(D:\mathcal X\to\D\) be a labelled operational data map built from normal
state/effect values, channels tested by specified preparations and effects,
represented modular or Osterwalder--Schrader data, DHR sector data, locally
covariant relative-Cauchy-evolution or probe-measurement data, relative entropy
or exact recovery data, crossed-product entropy data, or trace/area
assignments.

Assume that the finite central directions, when present, are read only through
one of the following: a normal quotient, a represented GNS quotient, a
central-state retraction, a diagonal subnet restriction, a generated algebra,
a specified nonseparating preparation/effect family, a restricted sector
functor, or a specified trace/area calibration.  Assume also that the
comparison class contains the corresponding finite central modifications unless
an explicit data exit or exclusion exit is imposed.

For each full-algebra target in the left column, \(D\) reconstructs that target
on the finite-central comparison class if and only if the matching exit in the
right column is supplied.
\begin{center}
\small
\begin{tabular}{>{\raggedright\arraybackslash}p{0.39\linewidth}
                >{\raggedright\arraybackslash}p{0.51\linewidth}}
\toprule
Target & Exact exit \\
\midrule
Full central support of the supplied normal readouts &
The readouts are faithful on the finite central positive cone. \\
Atoms and weights of a specified finite center &
The restrictions of the supplied central readouts separate the center. \\
Variable finite cardinality, including the two-atom predicate &
The data distinguish the corresponding cardinality fibers, or the
cardinality is supplied as calibrated data. \\
Factor-valuedness of a retained local algebra &
A factor, irreducibility, quotient, inadmissibility, or minimality
hypothesis excludes the nontrivial finite center. \\
A proposed spatial tensor product for a bipartition &
The tested pullbacks of multiplication have zero common kernel; with all
normal generated-algebra correlations this is injectivity of multiplication,
for example by split or W*-tensor-product independence. \\
Finite central channel dynamics &
Tested central preparations span the finite input algebra and tested central
effects separate the output algebra. \\
Central-character labels in sector data &
The supplied sector functor is faithful on the central-character fibers. \\
Absolute finite-center entropy, a horizon entropy-density coefficient, or the
Newton coupling inferred from \(S=A/(4G\hbar)\) &
The affine absolute calibration spans the claimed central state family,
the trace calibration spans the central coordinates used by the entropy term,
the geometric area scale is fixed, and every remaining same-data residual has
zero limiting area density; or an explicit quotient/exclusion removes the
residual. \\
\bottomrule
\end{tabular}
\end{center}
If the row exit is absent and the comparison class contains the corresponding
finite central modification, then there are \(X,Y\in\mathcal X\) with
\(D(X)=D(Y)\) and different values of the retained target.  Hence no
universal reconstruction rule from \(D\) can determine that target.
\end{theorem}

\begin{proof}
By Proposition \ref{prop:fiber}, every row is a fiber computation.

For support, restrict the supplied normal positive functionals \(F\) to the
finite center \(C\).  They rule out invisible central support exactly when
every nonzero central projection \(p\in C\) is detected by some
\(\varphi\in F\), i.e. \(\varphi(p)>0\).  If such a projection is missed, the
corner \(p\M p\) can be retained or adjoined without changing any supplied
readout.  If no projection is missed, equality of readouts forbids a hidden
positive central corner.  This is precisely central faithfulness.

For atoms and weights of a specified finite center
\(C\simeq\ell^\infty_n\), the readouts induce a linear map
\[
  r_F:C\to \C^F,\qquad r_F(c)=(\varphi(c))_{\varphi\in F}.
\]
The central coordinates are reconstructed exactly when \(r_F\) is injective,
or affinely injective on the normalized central state simplex.  If \(r_F\) is
not injective, two central observables or two central weights have identical
labelled pairings.  Across a comparison class in which \(n\) varies, the same
fiber criterion says that finite cardinality, and in particular the predicate
\(n=2\), is reconstructed only when the data distinguish the corresponding
cardinality fibers.

Factor-valuedness is the assertion that the relevant center is trivial.  A
finite direct sum of factor fibers has a nontrivial finite center unless only
one summand is present.  Thus any data fiber that leaves finite cardinality or
central separation undecided contains same-data candidates with different
factor-valuedness.  A factor, irreducibility, quotient, inadmissibility, or
minimality hypothesis removes that pair from the comparison class.

For tensor factorization, let
\(\A,\B\subset \N(O)\) be commuting von Neumann algebras.  Generated-algebra
correlations evaluate an element \(x\in\A\bar\otimes\B\) only through the
normal multiplication representation
\[
  \mu:\A\bar\otimes\B\to\A\vee\B,\qquad a\otimes b\mapsto ab .
\]
With tested functionals \(\Gamma\subset(\A\vee\B)_*\), two tensor elements
have the same generated-algebra data exactly when their difference lies in
\(\bigcap_{\gamma\in\Gamma}\ker(\gamma\circ\mu)\).  The unquotiented spatial
tensor coordinate therefore descends exactly when this common kernel is zero.
If all normal functionals on \(\A\vee\B\) are supplied, normal functionals
separate points of the generated algebra, so the common kernel is \(\ker\mu\).
Split or W*-tensor-product independence is the standard AQFT exit making this
multiplication representation faithful in the cases considered here
\cite{DoplicherLongoSplit,BuchholzWichmann,WernerSplit}.

For finite central process data, the channel table consists of pairings
\(e(K\rho)\) between tested preparations and tested effects.  The full finite
central channel is determined exactly when tested preparations span the finite
input vector space and tested effects separate the finite output vector space.
If either condition fails, a finite classical channel can be perturbed in an
untested direction without changing the displayed table.

For sector data, DHR/Doplicher--Roberts reconstruction applies to the sector
category actually supplied \cite{DHR1,DHR2,DRDuality,DR}.  If finite central
characters are forgotten by restriction to a diagonal subnet, those characters
are not present in the supplied category.  They are reconstructed exactly when
the sector functor remains faithful on the central-character fibers.

For entropy and horizon coefficients, finite central residuals are explicit.
Uniform \(q\)-fold multiplicity contributes \(\log q\) for one unresolved
finite center, or \(m(\Sigma)\log q\) for \(m(\Sigma)\) such centers on a cut
\(\Sigma\).  Finite trace weights contribute logarithmic terms of the form
\(\sum_i p_i\log\lambda_i\).  Relative-entropy and JLMS-type data determine
absolute entropy or area potentials only modulo affine functions on the tested
central state family.  Therefore an absolute entropy term or density
coefficient descends exactly when the supplied trace, state-family, affine,
and geometric-area calibrations make every same-data residual vanish on the
claimed family, or at least have zero limiting density after division by
area.  If area units and \(\hbar\) are fixed and \(c>0\), the map
\(c\mapsto (4\hbar c)^{-1}\) is injective, so the same criterion identifies
the Newton coupling inferred from \(S=A/(4G\hbar)\).

These row computations prove both directions.  If the exit holds, the target
is constant on finite-central fibers.  If it fails and the comparison class
contains the corresponding modification, one of the countermodels in the next
section gives a same-data pair with different target value.
\end{proof}

\paragraph{Closest prior results and novelty for Theorem \ref{thm:main}.}
The support part is basic von Neumann algebra central-support theory
\cite{Sakai,TakesakiI}.  The tensor positive exit is the split/W*-independence
literature \cite{DoplicherLongoSplit,BuchholzWichmann,WernerSplit}.  The sector
positive theorem is DHR/Doplicher--Roberts reconstruction from a supplied
category \cite{DHR1,DHR2,DRDuality,DR}.  The locally covariant and measurement
inputs are BFV and Fewster--Verch \cite{BFV,FewsterVerch,
FewsterVerchMeasurements}.  The entropy inputs are Araki/Petz/JLMS/OAQEC and
gravitational crossed products \cite{ArakiRelativeEntropy,Petz,JLMS,HarlowRT,
DHW,WittenCrossed,LongoWitten,CLPW,CPW}.  The theorem's claimed contribution is
the simultaneous finite-central iff: these positive frameworks reconstruct
full support, central cardinality, unquotiented tensor factors, retained sector
labels, and absolute finite-center entropy contributions exactly when the
corresponding finite-central exit is part of the supplied data or hypotheses.

\section{Three Countermodels}

\subsection{Quotient-visible support and finite rank}

Let \(\B\) and \(\M\) be nonzero von Neumann algebras.  Put
\[
  \widehat\B_0=\B,\qquad
  \widehat\B_q=\B\oplus(\M\bar\otimes\ell^\infty(\{1,\ldots,q\})),
  \quad q\ge2,
\]
and let \(\pi_q:\widehat\B_q\to\B\) be the normal quotient onto the first
summand, with \(\pi_0=\mathrm{id}_\B\).  Suppose every displayed datum is a
labelled function of \(\pi_q\), say
\[
  D_q=(\widetilde D_\lambda\circ\pi_q)_{\lambda\in\Lambda},
\]
with fixed labels and codomains.  Then the reported data are identical for
\(q=0,2,3\).  But for \(q\ge2\),
\[
  p_q=(0,1_{\M\bar\otimes\ell^\infty_q})
\]
is a nonzero central projection killed by every pulled-back readout, while no
such complementary killed projection exists for \(q=0\).  Thus quotient-visible
data do not determine full ambient support, faithfulness on \(\widehat\B_q\),
or the interpretation of the killed kernel as absent, null, gauge,
inadmissible, or physical.  Comparing \(q=2\) and \(q=3\) also shows that the
same data do not determine finite central cardinality or the two-atom
predicate.

This example is not a claim that the extra summand is physically present.  It
is a no-go for deriving its absence from data that factor through \(\pi_q\).
A physical theory evades the witness by supplying representation faithfulness,
full central tomography, a factor hypothesis, or an explicit quotient/gauge/null
or inadmissibility principle.

\subsection{Generated algebra versus spatial tensor product}

Let \(C=\ell^\infty(\{1,\ldots,n\})\) with \(n\ge2\), and view \(C\) twice as
commuting subalgebras of itself.  The generated algebra is again \(C\), while
the candidate spatial tensor product is \(C\bar\otimes C\).  Multiplication
acts by
\[
  \mu(e_i\otimes e_j)=\delta_{ij}e_i,
\]
where \(e_i\) are the minimal central projections.  Hence
\(e_i\otimes e_j\ne0\) for \(i\ne j\), but \(\mu(e_i\otimes e_j)=0\).  Normal
states on the generated algebra see only the diagonal image.  Normal product
states on \(C\bar\otimes C\) separate the off-diagonal tensor atom.

Thus generated-algebra tomography is not the same as spatial tensor-product
tomography.  The exact repair is the usual one: injectivity of multiplication,
for example from a split or W*-tensor-product independence hypothesis.

\subsection{Entropy density and Newton normalization}

Let the displayed entropy data be insensitive to an unresolved finite center on
a cut \(\Sigma\), but let an entropy convention count \(m(\Sigma)\) uniform
\(q\)-point centers.  The counted entropy changes by
\[
  \Delta S(\Sigma)=m(\Sigma)\log q .
\]
If \(A(\Sigma)\to\infty\) and
\[
  \frac{m(\Sigma)}{A(\Sigma)}\longrightarrow \beta ,
\]
then the entropy-density coefficient changes by \(\beta\log q\).  The same
displayed represented, retracted, relative-entropy, modular, or crossed-product
data can therefore correspond to two coefficients unless the finite-center
counting convention is calibrated, quotiented, or shown to have zero density.
After geometric area units and \(\hbar\) are fixed, a positive coefficient
\(c\) determines \(G=(4\hbar c)^{-1}\).  Thus a nonzero same-data residual in
\(c\) is also a nonidentifiability of the Newton coupling inferred from the
area law.

This is compatible with the positive Bekenstein--Hawking, JLMS/OAQEC, and
crossed-product entropy results \cite{Bekenstein,Hawking,JLMS,HarlowRT,DHW,
WittenCrossed,LongoWitten,CLPW,CPW}.  Those results supply or compute entropy
for specified algebras and normalizations.  The witness says only that a
represented or retracted data map that does not include the finite central
calibration cannot identify an absolute coefficient that counts it.

\section{Tensor Factorization and Horizon Coefficient}

\begin{theorem}[Tensor factors and entropy coefficient from operational data]
\label{thm:tensor-entropy}
Let \(O\mapsto\N(O)\) be a Haag--Kastler net on a fixed spacetime, or the
objectwise net of a locally covariant theory on a fixed globally hyperbolic
spacetime.  For each \(i\in\mathcal I\), let
\(\A_i,\B_i\subset\N(O_i)\) be commuting von Neumann algebras and let
\[
  \mu_i:\A_i\bar\otimes\B_i\to\A_i\vee\B_i,\qquad a\otimes b\mapsto ab
\]
be multiplication.  Let \(\Gamma_i\subset(\A_i\vee\B_i)_*\) be the supplied
normal correlation functionals.  Let \(E\) be a comparison class of entropy
assignments on cuts \(\Sigma\), with fixed geometric areas \(A(\Sigma)>0\),
\(A(\Sigma)\to\infty\), labelled entropy data \(D_E:E\to\D_E\), and coefficient
\[
  c(X)=\lim_{\Sigma}\frac{S_X(\Sigma)}{A(\Sigma)}
\]
when the limit exists.

For \(x_i\in\A_i\bar\otimes\B_i\), define the labelled data
\[
  D((x_i),X)=
  \left(
    (\gamma(\mu_i(x_i)))_{i\in\mathcal I,\gamma\in\Gamma_i},
    D_E(X)
  \right).
\]
The target is the unquotiented tensor family \((x_i)_{i\in\mathcal I}\)
together with \(c(X)\).  This target is reconstructible from \(D\) if and only
if
\[
  \bigcap_{\gamma\in\Gamma_i}\ker(\gamma\circ\mu_i)=0
  \quad\hbox{for every }i\in\mathcal I
\]
and
\[
  D_E(X)=D_E(Y)\Longrightarrow
  \lim_{\Sigma}\frac{S_X(\Sigma)-S_Y(\Sigma)}{A(\Sigma)}=0 .
\]
If \(\Gamma_i=(\A_i\vee\B_i)_*\), the first condition is \(\ker\mu_i=0\).  If
area units and \(\hbar\) are fixed and \(c>0\), the second condition is also the
criterion for reconstructing the Newton coupling
\[
  G=\frac{1}{4\hbar c}.
\]
\end{theorem}

\begin{proof}
Apply Proposition \ref{prop:fiber} coordinatewise.  For fixed \(i\), equality
of generated-algebra correlation data for \(x_i\) and \(y_i\) means
\[
  \gamma(\mu_i(x_i-y_i))=0\qquad(\gamma\in\Gamma_i),
\]
so \(x_i-y_i\) lies in the common tested multiplication kernel.  The
unquotiented tensor coordinate is constant on the correlation fiber exactly
when that kernel is zero.  If all normal functionals on the generated algebra
are supplied, normal functionals separate points, so the common kernel is
\(\ker\mu_i\).

The entropy coefficient is constant on \(D_E\)-fibers exactly when every
same-data residual has zero limiting density, which is the second displayed
condition.  The map \(c\mapsto(4\hbar c)^{-1}\) is injective for positive
\(c\), so the same criterion identifies the Newton coupling after the area
scale and \(\hbar\) are fixed.  Failure of either condition gives a same-data
pair by varying only the failed coordinate.
\end{proof}

\paragraph{Closest prior results and novelty for Theorem
\ref{thm:tensor-entropy}.}
The tensor positive direction is the standard split/W*-independence exit
\cite{DoplicherLongoSplit,BuchholzWichmann,WernerSplit}.  The entropy positive
inputs are Bekenstein--Hawking, FLM/JLMS/OAQEC, and Type II/crossed-product
entropy for specified algebras and normalizations
\cite{Bekenstein,Hawking,FLM,JLMS,HarlowRT,DHW,WittenCrossed,LongoWitten,CLPW,
CPW}.  The new packaging is the product fiber statement: generated-algebra
operational correlations and entropy data identify the joint target only when
the tensor-kernel condition and entropy-density calibration both hold.  Either
missing condition gives a same-data countermodel even if the other side is
perfectly calibrated.

\section{Relation to Positive Reconstruction Theorems}

\begin{longtable}{p{0.24\linewidth}p{0.31\linewidth}p{0.34\linewidth}}
\toprule
Positive input & What it supplies & Finite-central limitation isolated here \\
\midrule
Tomita--Takesaki theory and crossed products & Modular flow and crossed
product of a specified represented algebra and weight
\cite{TakesakiI,TakesakiII,TakesakiDuality,WittenCrossed}. & They do not
recover a finite central summand removed before representation; faithful full
centers are carried along only when the full algebra is supplied. \\
DHR/Doplicher--Roberts reconstruction & Field algebra and compact gauge group
from a supplied sector category \cite{DHR1,DHR2,DRDuality,DR}. & If diagonal
restriction forgets central characters, those characters are not in the input
category. \\
BFV/Fewster--Verch locally covariant QFT & Functorial dynamics, relative
Cauchy evolution, dynamical locality, and local measurement schemes for a
supplied functor and supplied probes
\cite{BFV,FewsterVerch,FewsterVerchMeasurements}. & Constant finite-center
copies remain same-data after retraction unless the probes or object algebras
separate the center. \\
Split inclusions & W*-tensor-product independence for a supplied bipartition
\cite{DoplicherLongoSplit,BuchholzWichmann,WernerSplit}. & This is exactly the
positive exit from the shared-center multiplication kernel; it is not implied
by generated-algebra tomography alone. \\
Process tomography & Channel reconstruction relative to supplied
preparations and effects \cite{Jamiolkowski,Choi,ChuangNielsenProcess}. &
Untested finite central input or output directions remain invisible. \\
No-broadcasting & Noncommuting mixed states cannot be broadcast by a universal
physical channel \cite{NoBroadcasting}. & The present obstruction is different:
finite centers are commutative once supplied, but operational data that omit a
central coordinate cannot infer which central coordinate or cardinality was
present. \\
JLMS/OAQEC and crossed-product entropy & Entropy, recovery, or Type II
algebra statements for specified algebras and normalizations
\cite{Petz,JLMS,HarlowRT,DHW,WittenCrossed,LongoWitten,CLPW,CPW}. & Absolute
finite-center entropy terms and area coefficients require trace/area
calibration or zero-density residuals. \\
\bottomrule
\end{longtable}

The table is meant to prevent an inflated novelty claim.  The note does not
show that finite centers occur in nature and does not refute the positive
theorems.  It gives a test for proposed operational reconstructions: either
the target factors through the supplied data, the hypotheses exclude the
finite-central comparison pair, or the data contain the exact separator or
calibration.  If none of these is true, the three countermodels above supply a
same-data pair with different target.

\section{Assumption Ledger for Referees}

This section is included to make the submission checkable without reading the
longer manuscript.  Each row separates a standard cited input from the new
finite-central descent step and from the physical input that is deliberately
not proved here.

\begin{longtable}{>{\raggedright\arraybackslash}p{0.26\linewidth}
                  >{\raggedright\arraybackslash}p{0.32\linewidth}
                  >{\raggedright\arraybackslash}p{0.30\linewidth}}
\toprule
Hypothesis used & Why it is standard & What this note does not prove \\
\midrule
Normal quotients and represented support summands of von Neumann algebras. &
Weak-* closed ideals are central summands, and normal representations have
central support kernels \cite{Sakai,TakesakiI}. & That a proposed physical
theory admits no additional central summand, or that every quotient-zero
direction is physically null. \\
Normal readouts restricted to a finite abelian center. & Finite-dimensional
linear algebra: coordinate functionals separate
\(\ell^\infty_n\), and positive functionals detect support exactly when every
central atom is seen. & That a given operational principle supplies all such
central readouts as physical effects. \\
Generated-algebra correlations for commuting subalgebras. & Correlations on
\(\A\vee\B\) are functions of the multiplication image
\(\mu:\A\bar\otimes\B\to\A\vee\B\). & That \(\mu\) is injective; this is the
split/W*-independence exit when it holds. \\
Sector reconstruction from a supplied category. & DHR/Doplicher--Roberts
reconstruction starts from the retained category of sectors
\cite{DHR1,DHR2,DRDuality,DR}. & That a diagonal restriction has retained
central-character labels it has already forgotten. \\
Locally covariant relative Cauchy evolution and local probes. & BFV and
Fewster--Verch formulate these notions for a supplied functor and supplied
probe couplings \cite{BFV,FewsterVerch,FewsterVerchMeasurements}. & That every
finite central direction has a source or induced probe effect separating it. \\
Process tomography. & Choi--Jamiolkowski tomography reconstructs a channel
relative to an informationally complete input/effect domain
\cite{Jamiolkowski,Choi,ChuangNielsenProcess}. & That the admitted physical
preparations and effects include the missing finite central directions. \\
Relative entropy, exact recovery, and crossed-product entropy. & These are
positive results for specified algebras, channels, traces, and state families
\cite{ArakiRelativeEntropy,Petz,JLMS,HarlowRT,DHW,WittenCrossed,LongoWitten,
CLPW,CPW}. & That an absolute finite-center entropy term or geometric area
scale has been calibrated when the data include only relative or represented
quantities. \\
\bottomrule
\end{longtable}

The last column is not a defect in the cited theorems.  It is the boundary of
this note.  A theory of quantum gravity may have a compelling physical reason
for any of these exits: factor-valued local algebras, a gauge quotient, a
superselection analysis retaining the relevant central characters, a split
inclusion, or an absolute trace/area normalization.  The point is that such a
reason must appear as an assumption or theorem.  It is not generated by data
that were, by construction, constant on the finite-central comparison pair.

\section{Novelty Boundary by Claim}

\paragraph{Support and finite cardinality.}
The closest known ingredients are standard central support and quotient facts
for von Neumann algebras.  The new claim is not that central support exists, but
that a quotient-visible operational package cannot reconstruct full ambient
support, the kernel interpretation, or finite central rank unless the package
contains a central separator or an exclusion rule.  The \(q=0\) versus
\(q=2\) comparison tests support totality; the \(q=2\) versus \(q=3\)
comparison tests finite central cardinality and the two-atom predicate.

\paragraph{Tensor factorization.}
The closest positive result is the split property and W*-tensor-product
independence.  The note's contribution is the converse fiber formulation for
operational data: generated-algebra correlations reconstruct only the
multiplication image, so the unquotiented spatial tensor coordinate is
identified exactly when the tested multiplication kernel vanishes.  This is a
useful warning for operational reconstructions because full tomography of
\(\A\vee\B\) is still tomography of the quotient by \(\ker\mu\), not of an
arbitrary proposed tensor product.

\paragraph{Sectors and locally covariant data.}
DHR/DR and BFV/Fewster--Verch are positive frameworks for supplied data.  The
finite-center point is that restriction, retraction, or copied-center
amplification can make distinct central-character or finite-center theories
look identical to those supplied data.  If a reconstruction theorem claims the
full center anyway, it must include a sector-character separator, a probe
effect/source separating the center, or an exclusion principle forbidding the
comparison pair.

\paragraph{Broadcasting and copying.}
No-broadcasting theorems rule out a universal physical process broadcasting
noncommuting mixed states.  This note does not strengthen that theorem.  Its
finite centers are classical once included in the algebra; the obstruction is
that represented, retracted, or restricted data may have omitted the central
coordinate before any copying question is posed.  The repair is therefore
central separation or an exclusion principle, not a broadcasting channel.

\paragraph{Horizon entropy coefficient.}
Bekenstein--Hawking, FLM/JLMS, OAQEC, and gravitational crossed-product papers
provide positive entropy formulas for specified geometric and algebraic
normalizations.  The finite-center descent statement is narrower: relative
entropy, exact recovery, or a represented Type II algebra does not by itself
identify an absolute finite-center counterterm if an affine, trace-weight, or
counting residual has nonzero area density and is invisible to the data.  Thus
the exact one-quarter coefficient requires an identifiable calibration or a
proof that every same-data residual has zero density.

\paragraph{Lean core.}
The Lean file is not a formal proof of AQFT.  Its novelty status is modest but
useful: it mechanically checks the finite logical fibers, quotient-factored
no-go lemmas, diagonal tensor-kernel witness, and entropy-density descent
criterion.  This removes the risk that the finite part of the argument is a
verbal sleight of hand.  The infinite-dimensional reductions remain standard
mathematical background, cited rather than formalized.

\section{What Was Cut From the Long Manuscript}

The longer manuscript contains many variants: modular-core refinements, OS
specializations, process-tomography examples, local-probe protocols,
sector-character wrappers, affine/JLMS calibration variants, and several Lean
wrappers.  They are omitted here unless they are needed for one of the two main
statements or for the three countermodels.  This is intentional.  A referee for
the short note should have to check only the following chain:
\begin{enumerate}
\item the fiber criterion in Proposition \ref{prop:fiber};
\item the standard operator-algebra facts cited in the assumption ledger;
\item the row-by-row finite-center calculation in Theorem \ref{thm:main};
\item the three same-data witnesses;
\item the coordinatewise tensor/entropy-density criterion in Theorem
      \ref{thm:tensor-entropy};
\item the Lean file as a finite logical sanity check, not as evidence for any
      analytic AQFT premise.
\end{enumerate}
Everything else is deferred to the longer manuscript or to future work.  In
particular, this note does not attempt to prove a microscopic origin of a
finite center, does not prove that physical local algebras fail to be factors,
does not derive the Bekenstein--Hawking coefficient, and does not classify all
locally covariant theories.  It gives a finite-central obstruction that a
positive reconstruction theorem must evade explicitly.

\section{Lean Core}

The accompanying Lean 4 file
\[
  \texttt{formalization/FiniteCenterSupport.lean}
\]
does not formalize AQFT, Tomita--Takesaki theory, DHR sectors, or crossed
products.  It formalizes the finite logical kernels that remain after the
standard analytic hypotheses reduce the question to finite central fibers.
The most relevant named statements are:
\[
\begin{gathered}
\texttt{quotientFactoredInvariantNoGoCore},\\
\texttt{quotientFactoredPropertyNoGoCore},\\
\texttt{quotientVisibleUniversalNoGoCore},\\
\texttt{finiteCenterOperationalCompletenessCore},\\
\texttt{finiteCenterLocalNetClassificationCore},\\
\texttt{localNetFamilyEntropyDescentCore},\\
\texttt{refereeTensorEntropyCriterionCore},\\
\texttt{universalFiniteCentralWitnessesCore}.
\end{gathered}
\]
These lemmas mechanically check the fiber criterion, quotient-factored
nonreconstruction, finite support coverage, finite coordinate separation,
diagonal multiplication kernels, finite process-table completeness, and
entropy-density descent.  The file therefore verifies the finite part of the
paper's logic while leaving the infinite-dimensional operator-algebraic inputs
to the cited literature.

\section{Conclusion}

The finite-central obstruction is small but sharp.  Operational data that have
already passed through a quotient, representation, retraction, generated
algebra, restricted test family, or incomplete calibration cannot reconstruct a
full-algebra target that changes inside the corresponding data fiber.  For
emergent-spacetime reconstructions, this means that tensor factorization and
the exact horizon entropy coefficient require identifiable extra hypotheses:
split/tensor independence for the former, and trace/area/entropy calibration
with zero same-data density residual for the latter.  Those hypotheses may be
physically correct in a given theory.  The theorem says only that they must be
present and checkable; they cannot be replaced by operational data that have
already forgotten the finite center.

\begin{thebibliography}{99}
\raggedright

\bibitem{Sakai}
S. Sakai, \emph{\(C^*\)-Algebras and \(W^*\)-Algebras}, Springer, 1971.

\bibitem{TakesakiI}
M. Takesaki, \emph{Theory of Operator Algebras I}, Springer, 1979.

\bibitem{TakesakiII}
M. Takesaki, \emph{Theory of Operator Algebras II}, Springer, 2003.

\bibitem{TakesakiDuality}
M. Takesaki,
``Duality for crossed products and the structure of von Neumann algebras of
type III,'' Acta Math. \textbf{131} (1973) 249--310,
doi:10.1007/BF02392041.

\bibitem{OS1}
K. Osterwalder and R. Schrader,
``Axioms for Euclidean Green's functions,''
Commun. Math. Phys. \textbf{31} (1973) 83--112,
doi:10.1007/BF01645738.

\bibitem{OS2}
K. Osterwalder and R. Schrader,
``Axioms for Euclidean Green's functions II,''
Commun. Math. Phys. \textbf{42} (1975) 281--305,
doi:10.1007/BF01608978.

\bibitem{HaagKastler}
R. Haag and D. Kastler,
``An algebraic approach to quantum field theory,''
J. Math. Phys. \textbf{5} (1964) 848--861,
doi:10.1063/1.1704187.

\bibitem{BFV}
R. Brunetti, K. Fredenhagen and R. Verch,
``The generally covariant locality principle -- A new paradigm for local
quantum physics,''
Commun. Math. Phys. \textbf{237} (2003) 31--68,
\href{https://arxiv.org/abs/math-ph/0112041}{arXiv:math-ph/0112041},
doi:10.1007/s00220-003-0815-7.

\bibitem{FewsterVerch}
C. J. Fewster and R. Verch,
``Dynamical locality and covariance: What makes a physical theory the same in
all spacetimes?''
Ann. Henri Poincare \textbf{13} (2012) 1613--1674,
\href{https://arxiv.org/abs/1106.4785}{arXiv:1106.4785},
doi:10.1007/s00023-012-0165-0.

\bibitem{FewsterVerchMeasurements}
C. J. Fewster and R. Verch,
``Quantum fields and local measurements,''
Commun. Math. Phys. \textbf{378} (2020) 851--889,
\href{https://arxiv.org/abs/1810.06512}{arXiv:1810.06512},
doi:10.1007/s00220-020-03800-6.

\bibitem{DHR1}
S. Doplicher, R. Haag and J. E. Roberts,
``Fields, observables and gauge transformations I,''
Commun. Math. Phys. \textbf{13} (1969) 1--23,
doi:10.1007/BF01645267.

\bibitem{DHR2}
S. Doplicher, R. Haag and J. E. Roberts,
``Local observables and particle statistics I,''
Commun. Math. Phys. \textbf{23} (1971) 199--230,
doi:10.1007/BF01877742.

\bibitem{DRDuality}
S. Doplicher and J. E. Roberts,
``A new duality theory for compact groups,''
Invent. Math. \textbf{98} (1989) 157--218,
doi:10.1007/BF01388849.

\bibitem{DR}
S. Doplicher and J. E. Roberts,
``Why there is a field algebra with a compact gauge group describing the
superselection structure in particle physics,''
Commun. Math. Phys. \textbf{131} (1990) 51--107,
doi:10.1007/BF02097680.

\bibitem{DoplicherLongoSplit}
S. Doplicher and R. Longo,
``Standard and split inclusions of von Neumann algebras,''
Invent. Math. \textbf{75} (1984) 493--536,
doi:10.1007/BF01388641.

\bibitem{BuchholzWichmann}
D. Buchholz and E. H. Wichmann,
``Causal independence and the energy-level density of states in local quantum
field theory,'' Commun. Math. Phys. \textbf{106} (1986) 321--344,
doi:10.1007/BF01454978.

\bibitem{WernerSplit}
R. F. Werner,
``Local preparability of states and the split property in quantum field
theory,'' Lett. Math. Phys. \textbf{13} (1987) 325--329,
doi:10.1007/BF00401161.

\bibitem{Jamiolkowski}
A. Jamio\l{}kowski,
``Linear transformations which preserve trace and positive semidefiniteness of
operators,'' Rep. Math. Phys. \textbf{3} (1972) 275--278,
doi:10.1016/0034-4877(72)90011-0.

\bibitem{Choi}
M.-D. Choi,
``Completely positive linear maps on complex matrices,''
Linear Algebra Appl. \textbf{10} (1975) 285--290,
doi:10.1016/0024-3795(75)90075-0.

\bibitem{ChuangNielsenProcess}
I. L. Chuang and M. A. Nielsen,
``Prescription for experimental determination of the dynamics of a quantum
black box,'' J. Mod. Opt. \textbf{44} (1997) 2455--2467,
\href{https://arxiv.org/abs/quant-ph/9610001}{arXiv:quant-ph/9610001},
doi:10.1080/09500349708231894.

\bibitem{NoBroadcasting}
H. Barnum, C. M. Caves, C. A. Fuchs, R. Jozsa and B. Schumacher,
``Noncommuting mixed states cannot be broadcast,''
Phys. Rev. Lett. \textbf{76} (1996) 2818--2821,
\href{https://arxiv.org/abs/quant-ph/9511010}{arXiv:quant-ph/9511010},
doi:10.1103/PhysRevLett.76.2818.

\bibitem{ArakiRelativeEntropy}
H. Araki,
``Relative entropy of states of von Neumann algebras,''
Publ. Res. Inst. Math. Sci. \textbf{11} (1976) 809--833,
doi:10.2977/prims/1195191148.

\bibitem{Petz}
D. Petz,
``Sufficient subalgebras and the relative entropy of states of a von Neumann
algebra,'' Commun. Math. Phys. \textbf{105} (1986) 123--131,
doi:10.1007/BF01212345.

\bibitem{Bekenstein}
J. D. Bekenstein,
``Black holes and entropy,'' Phys. Rev. D \textbf{7} (1973) 2333--2346,
doi:10.1103/PhysRevD.7.2333.

\bibitem{Hawking}
S. W. Hawking,
``Particle creation by black holes,''
Commun. Math. Phys. \textbf{43} (1975) 199--220,
doi:10.1007/BF02345020.

\bibitem{FLM}
T. Faulkner, A. Lewkowycz and J. Maldacena,
``Quantum corrections to holographic entanglement entropy,''
JHEP \textbf{11} (2013) 074,
\href{https://arxiv.org/abs/1307.2892}{arXiv:1307.2892},
doi:10.1007/JHEP11(2013)074.

\bibitem{JLMS}
D. L. Jafferis, A. Lewkowycz, J. Maldacena and S. J. Suh,
``Relative entropy equals bulk relative entropy,''
JHEP \textbf{06} (2016) 004,
\href{https://arxiv.org/abs/1512.06431}{arXiv:1512.06431},
doi:10.1007/JHEP06(2016)004.

\bibitem{HarlowRT}
D. Harlow,
``The Ryu--Takayanagi formula from quantum error correction,''
Commun. Math. Phys. \textbf{354} (2017) 865--912,
\href{https://arxiv.org/abs/1607.03901}{arXiv:1607.03901},
doi:10.1007/s00220-017-2904-z.

\bibitem{DHW}
X. Dong, D. Harlow and A. C. Wall,
``Reconstruction of bulk operators within the entanglement wedge in
gauge-gravity duality,'' Phys. Rev. Lett. \textbf{117} (2016) 021601,
\href{https://arxiv.org/abs/1601.05416}{arXiv:1601.05416},
doi:10.1103/PhysRevLett.117.021601.

\bibitem{WittenCrossed}
E. Witten,
``Gravity and the crossed product,''
JHEP \textbf{10} (2022) 008,
\href{https://arxiv.org/abs/2112.12828}{arXiv:2112.12828},
doi:10.1007/JHEP10(2022)008.

\bibitem{LongoWitten}
R. Longo and E. Witten,
``A note on continuous entropy,''
\href{https://arxiv.org/abs/2202.03357}{arXiv:2202.03357}.

\bibitem{CLPW}
V. Chandrasekaran, R. Longo, G. Penington and E. Witten,
``An algebra of observables for de Sitter space,''
JHEP \textbf{02} (2023) 082,
\href{https://arxiv.org/abs/2206.10780}{arXiv:2206.10780},
doi:10.1007/JHEP02(2023)082.

\bibitem{CPW}
V. Chandrasekaran, G. Penington and E. Witten,
``Large \(N\) algebras and generalized entropy,''
JHEP \textbf{04} (2023) 009,
\href{https://arxiv.org/abs/2209.10454}{arXiv:2209.10454},
doi:10.1007/JHEP04(2023)009.

\end{thebibliography}

\end{document}
